\documentclass[conference,a4paper]{APSIPA2026}
\usepackage{amsmath,amssymb}
\usepackage{graphicx}
\usepackage{multirow}
\usepackage{threeparttable}
\usepackage{booktabs}
\usepackage{tikz}
\usetikzlibrary{positioning,arrows.meta,calc,shapes.geometric,shapes.misc,fit,backgrounds}
\usepackage[backend=biber,style=ieee]{biblatex}
\addbibresource{refs.bib}

\usepackage{geometry}
\geometry{a4paper, top=19mm, bottom=43mm, right=13mm, left=13mm}

\usepackage{fancyhdr}
\renewcommand{\headrulewidth}{0pt}
\renewcommand{\footrulewidth}{0pt}
\fancypagestyle{firststyle}{
  \fancyhf{}
  \fancyhead[C]{2026 Asia Pacific Signal and Information Processing Association Annual Summit and Conference (APSIPA ASC)}
}

\begin{document}

\title{Multimodal Mamba Fusion for Noise-Robust\\
EEG Biometric Identification}

\author{
\authorblockN{
Ly Ngoc Vu\authorrefmark{1}\authorrefmark{2},
Huynh Cong Bang\authorrefmark{1}\authorrefmark{2} and
Nguyen Minh Hanh\authorrefmark{1}\authorrefmark{3}
}

\authorblockA{
\authorrefmark{1}
Faculty of Information Technology, Industrial University of Ho Chi Minh City, Vietnam \\
E-mail: dezzhuge@gmail.com, bang\_01036047@iuh.edu.vn, hanh\_01036046@iuh.edu.vn \\
\authorrefmark{2}Equal contribution (co-first authors). \quad
\authorrefmark{3}Corresponding author.}
}

\maketitle
\thispagestyle{firststyle}
\pagestyle{empty}

\begin{abstract}
EEG biometric identification is brittle under deployment noise and
unseen-subject conditions. Extending our two-stage WaveNet+ResNet
pipeline, we propose a multimodal Mamba framework: a Mamba-augmented
denoiser, a parallel spectrogram branch with frequency- and time-axis
Mamba sweeps, and a gated self-attention fusion module with auxiliary
branch losses. On the PhysioNet AEP dataset under the AEP-Hybrid
protocol (16 known + 4 held-out impostor subjects, three seeds,
inheriting the partition of~\cite{vu2026aisei}), the full model (V4)
attains the best
Precision@1 on all three noise families ($82.4\%$ Gaussian,
$87.7\%$ power-line, $83.8\%$ EMG) with Stage-I SI-SNR
$12.15$/$32.38$/$13.83$\,dB. Held-out-subject AUROC $\approx 0.50$
motivates calibration as future work.
\end{abstract}

\begin{IEEEkeywords}
EEG biometrics, denoising autoencoder, state-space models, Mamba,
multimodal fusion, gated self-attention fusion, metric learning.
\end{IEEEkeywords}

\section{Introduction}
The push toward unobtrusive authentication for wearables, brain-computer
interfaces, and cyber-physical systems has revived interest in
EEG-based biometrics: the signal is internal, hard to spoof, and rich in
subject-specific micro-patterns~\cite{fallahi2023brainnet,yang2022eeg,zhang2017mindid}.
Concretely, we treat EEG biometrics as the following task: given a
4-channel, $T=800$-sample EEG window from an enrolled subject, possibly
corrupted by deployment noise, output a unit-norm $128$-d biometric
embedding such that (a) the closest enrolled identity in cosine
distance is the correct subject (closed-set retrieval), and (b) an
unseen subject can in principle be flagged via a distance threshold
(open-set verification). Two practical challenges persist:
real recordings are corrupted by power-line interference, muscle
artifacts, and broadband noise~\cite{alzahab2024aep}, and deployable
systems must generalize to subjects that were never seen during
training~\cite{deng2019arcface,wang2019multisim}.

Our prior two-stage pipeline~\cite{vu2026aisei} decouples a WaveNet
denoiser from a ResNet metric-learning embedder but operates on a
single view of the signal with a fixed convolutional receptive field.
This paper grounds three contributions in concrete observations on
that design.
\textbf{(i)} The denoiser's receptive field is fixed by depth and
dilation, so we insert a \emph{selective state-space}
(Mamba)~\cite{gu2023mamba,gu2022s4} block at the midpoint of the
WaveNet stack for content-aware mixing in linear time.
\textbf{(ii)} The raw waveform hides rhythmic line-noise structure, so
we add a \emph{spectrogram branch} that re-encodes the denoised signal
into a 2-D time-frequency map with frequency- and time-axis Mamba
sweeps.
\textbf{(iii)} A single embedding view risks branch-specific bias, so
we fuse the two embeddings with a \emph{gated self-attention fusion}
module, jointly trained with auxiliary metric losses on each branch.

We evaluate four progressively richer architectures on the PhysioNet
Auditory-Evoked-Potential EEG dataset~\cite{alzahab2024aep}: V1 is the
WaveNet-only denoiser baseline of~\cite{vu2026aisei}; V2 inserts a
Mamba block at the WaveNet midpoint; V3 is a single-seed H100 quick
run with a tuned Mamba preset, reported as a directional reference;
V4 is the full multimodal proposal. Under the AEP-Hybrid protocol
(Sec.~\ref{sec:setup}), V4 delivers the highest mean closed-set
Precision@1 on every noise family in our benchmark and reaches up to
$32.4$\,dB SI-SNR on power-line noise.

\section{Related Work}
\textbf{EEG biometrics.}
Early EEG biometric systems relied on hand-crafted spectral
features~\cite{fallahi2023brainnet}. Deep models then took
over, including attention-recurrent
architectures~\cite{zhang2017mindid}, triplet-loss embeddings on raw
signals~\cite{yang2022eeg}, and Siamese networks for cross-session
matching~\cite{fallahi2023brainnet}. Compact convolutional networks such
as EEGNet~\cite{lawhern2018eegnet} and deep ConvNets for EEG
decoding~\cite{schirrmeister2017deep} provide widely used backbones, but
they are generally trained for closed-set classification and do not
explicitly target open-set verification.

\textbf{Denoising and signal enhancement.}
WaveNet-style dilated convolutions are a strong, well-studied tool for
raw-waveform denoising~\cite{vandenoord2016wavenet,luo2019tasnet}. Our
prior work~\cite{vu2026aisei} adopts this denoiser with an SI-SNR
objective; the current paper investigates whether replacing part of the
convolutional pathway with a selective state-space block produces a
qualitatively cleaner signal.

\textbf{State-space models.}
Structured state-space sequence models (S4)~\cite{gu2022s4} and the more
recent Mamba~\cite{gu2023mamba} approximate long-range temporal
dependencies in linear time, with content-dependent gating. They are
attractive for biomedical time-series where Transformer attention is
costly~\cite{vaswani2017attention} and where physiological non-stationarity
makes hand-tuned filters fragile. EEG-specific Mamba variants have begun
to appear for cross-subject decoding~\cite{wang2024eegmamba}; to our
knowledge, ours is the first study that combines a Mamba-augmented
denoiser with a Mamba spectrogram branch and gated self-attention fusion for
EEG biometric identification with open-set verification analysis.

\textbf{Metric learning and multimodal fusion.}
Angular margin and pair-weighted losses (ArcFace~\cite{deng2019arcface},
multi-similarity~\cite{wang2019multisim}, contrastive~\cite{hadsell2006contrastive},
triplet~\cite{schroff2015facenet}) are standard tools for open-set
verification. Audio spectrogram transformers~\cite{gong2021ast}
demonstrate that converting a 1-D acoustic signal into a 2-D
time-frequency map and applying sequence models on top is effective; we
adopt the same intuition with Mamba sweeps along both spectrogram axes
and use a gated self-attention fusion to combine the two views, instead
of a simpler concatenation or score-level fusion.

\section{Method}
\subsection{Overall Pipeline}
The pipeline keeps the two-stage decomposition introduced in
\cite{vu2026aisei} but enriches each stage. Stage~I is responsible for
removing deployment noise from the raw waveform, and Stage~II is
responsible for mapping the cleaned signal to a 128-d biometric
embedding. Stage~I is a WaveNet+Mamba denoiser trained with a
scale-invariant signal-to-noise ratio (SI-SNR) loss; Stage~II is a
multimodal embedder that consumes the denoised waveform together with a
short-time Fourier transform (STFT)-based spectrogram and emits a
128-d L2-normalized biometric embedding (the L2-normalization layer
is what produces the unit-norm vectors used downstream). The two stages are trained
sequentially: Stage~I for $30$ epochs and then frozen as the upstream
module for Stage~II.

Fig.~\ref{fig:v4-arch} details the V4 multimodal architecture, which
augments the second stage with a spectrogram branch and gated
self-attention fusion.

\subsection{Stage I: noise-conditioned waveform denoiser}
The denoiser ingests a 4-channel raw EEG epoch
$x_{\mathrm{noisy}}\in\mathbb{R}^{4\times 800}$ and emits a clean estimate
$\hat{x}=\mathcal{D}(x_{\mathrm{noisy}})$. The WaveNet backbone
contributes multiscale dilated convolutions for raw-waveform modeling,
while the inserted Mamba block contributes selective, content-aware
long-range mixing in $\mathcal{O}(L)$ operations. We retain the
dilated-convolution backbone of
\cite{vandenoord2016wavenet,luo2019tasnet}: three blocks of four
gated dilated layers with dilation $2^{i}$ for $i\in\{0,1,2,3\}$. At the
midpoint of the stack we insert a single \emph{Mamba block} that performs
content-dependent state-space mixing over the temporal axis with
$d_{\mathrm{model}}=64$, $d_{\mathrm{state}}=16$, kernel $4$, and
expand factor $2$. A pre-LayerNorm and a residual connection wrap the
block. Architecturally, this swaps a single fixed-receptive-field
convolutional layer for a selective state-space layer that re-weights
information across the entire epoch in $\mathcal{O}(L)$ operations.

We synthesize paired training data, i.e. (noisy input, clean target)
EEG pairs, by mixing clean EEG with three noise families: Gaussian,
$50$\,Hz power-line, and broadband EMG-like ($20$--$80$\,Hz) at
controlled SNR levels. The training objective is the SI-SNR loss, a
scale-invariant reconstruction-quality indicator standard in source
separation~\cite{luo2019tasnet}, used here so that Stage~I is robust
to amplitude rescaling between the clean target and the reconstructed
estimate. The denoiser is trained for $30$ epochs:
\begin{equation}
\begin{aligned}
\mathcal{L}_{\mathrm{SI\text{-}SNR}}(\hat{x},x_{\mathrm{clean}})
&= -10\,\log_{10}\!\frac{\|\alpha x_{\mathrm{clean}}\|^{2}}
{\|\hat{x}-\alpha x_{\mathrm{clean}}\|^{2}},\\
\alpha &= \frac{\langle \hat{x}, x_{\mathrm{clean}}\rangle}
{\|x_{\mathrm{clean}}\|^{2}}.
\end{aligned}
\end{equation}

\subsection{Stage II: dual-view biometric embedder}
\textbf{EEG branch.} The denoised epoch is reshaped from
$(4,800)$ to $(4,25,32)$ and passed through a ResNet-18 or ResNet-34
backbone~\cite{he2016resnet} with a $3{\times}3$ stride-1 stem and no
maxpool to preserve the small spatial structure, followed by a projection
head $\text{Linear}(\cdot,256)\to\text{ReLU}\to\text{Dropout}(0.1)
\to\text{Linear}(256,128)\to\text{BatchNorm}$ and L2 normalization,
producing the unit-norm waveform embedding $e_{\mathrm{eeg}}\in\mathbb{R}^{128}$.

\textbf{Spectrogram branch.} We compute a short-time Fourier transform
of the denoised signal with $n_{\mathrm{fft}}=128$ and hop length~$64$,
yielding a 4-channel time-frequency tensor $S\in\mathbb{R}^{4\times 65\times 13}$.
A two-layer convolutional stem ($4{\to}32{\to}64$ channels with $3{\times}3$
kernels and BatchNorm--ReLU) projects $S$ into a 64-d feature map.
Following the audio-spectrogram-transformer
intuition~\cite{gong2021ast} of decomposing a 2-D map into two 1-D
sequence problems, we then apply two Mamba sweeps with a shared
hidden width of~$64$: \emph{frequency-axis} (sequence is the
frequency axis with the time axis flattened into the batch) and
\emph{time-axis} (sequence is the time axis with frequency
flattened); the first sweep models per-band rhythmic structure, the
second models its temporal evolution. After mean-pooling over the spatial
dimensions, an MLP head produces the unit-norm spectrogram embedding
$e_{\mathrm{spec}}\in\mathbb{R}^{128}$ via a final L2-normalization
layer that is what makes the embedding unit-norm.

\textbf{Gated self-attention fusion.} The two unit-norm vectors are
stacked into a length-2 token sequence
$T=[e_{\mathrm{eeg}};e_{\mathrm{spec}}]\in\mathbb{R}^{2\times 128}$ and
passed through a four-head self-attention block (shared $Q/K/V$
derived from $T$; referred to as a cross-attention module
in~\cite{vu2026aisei}). The attended tokens
$\tilde{T}=[\tilde{e}_{\mathrm{eeg}};\tilde{e}_{\mathrm{spec}}]$ are
mixed by a learned gate
$g=\sigma(W_g[e_{\mathrm{eeg}};e_{\mathrm{spec}}])\in(0,1)^{128}$
computed from the pre-attention embeddings, where $\sigma$ is the
elementwise sigmoid and $W_g\in\mathbb{R}^{128\times 256}$ is a
learned linear projection acting as a soft per-coordinate selector
between the two views,
\begin{equation}
e_{\mathrm{mix}}=g\odot \tilde{e}_{\mathrm{eeg}}
+(1-g)\odot \tilde{e}_{\mathrm{spec}},
\end{equation}
followed by a residual MLP and L2 normalization to give the fused
biometric embedding $f\in\mathbb{R}^{128}$ used at inference.

\textbf{Joint training.} Stage~II is trained for 30 epochs with one of
two metric losses on $f$: ArcFace~\cite{deng2019arcface} or
multi-similarity~\cite{wang2019multisim}, drawn from
\cite{musgrave2020pml}. To prevent the fusion module from collapsing onto
one branch, we add weighted auxiliary metric losses on each branch:
\begin{equation}
\mathcal{L}_{\mathrm{II}}=\mathcal{L}_{\mathrm{metric}}(f)
+\lambda_{e}\mathcal{L}_{\mathrm{metric}}(e_{\mathrm{eeg}})
+\lambda_{s}\mathcal{L}_{\mathrm{metric}}(e_{\mathrm{spec}}).
\end{equation}
Numerical values for $\lambda_{e}$, $\lambda_{s}$, and the
\texttt{MPerClassSampler} parameter $m$ are fixed across all
experiments and reported in Sec.~\ref{sec:setup}; the chosen
operating point lies on the plateau of single-seed sensitivity
sweeps that we defer to the journal version.

\begin{figure*}[t]
\centering
\resizebox{\textwidth}{!}{%% v4_arch.tex — V4 multimodal architecture diagram (compact horizontal layout)
%% Rendered via \resizebox{\textwidth}{!}{\input{...}} inside figure*

\usetikzlibrary{positioning, arrows.meta, calc, fit, backgrounds, shapes.geometric, shapes.misc}

%% ── Colour definitions (muted academic palette) ──────────────────────────────
\definecolor{S1Blue}{RGB}{52,100,158}
\definecolor{S1BgCol}{RGB}{232,240,252}
\definecolor{S1Bdr}{RGB}{110,150,200}

\definecolor{BranchACol}{RGB}{180,100,30}
\definecolor{BranchABg}{RGB}{253,245,228}
\definecolor{BranchABdr}{RGB}{200,140,60}

\definecolor{BranchBCol}{RGB}{38,120,90}
\definecolor{BranchBBg}{RGB}{228,248,238}
\definecolor{BranchBBdr}{RGB}{70,160,110}

\definecolor{FuseCol}{RGB}{130,30,50}
\definecolor{FuseBg}{RGB}{250,232,236}
\definecolor{FuseBdr}{RGB}{175,60,80}

\definecolor{IOCol}{RGB}{60,60,70}
\definecolor{IOBg}{RGB}{240,240,244}

\definecolor{AuxCol}{RGB}{185,40,40}
\definecolor{AuxBg}{RGB}{255,238,238}

\definecolor{ShapeCol}{RGB}{80,80,95}
\definecolor{StageLbl}{RGB}{40,40,55}

\begin{tikzpicture}[
    font=\small\sffamily,
    node distance=5mm and 7mm,
    >={Stealth[length=2.0mm, width=1.4mm]},
    every node/.style={align=center},
    blk/.style={
        draw, rounded corners=2pt,
        inner sep=4pt, minimum height=11mm,
        line width=0.55pt
    },
    ionode/.style={
        draw=IOCol, fill=IOBg, rounded corners=1pt,
        inner sep=3pt, minimum height=9mm, minimum width=16mm,
        line width=0.5pt, font=\small\sffamily
    },
    s1blk/.style={blk, draw=S1Blue, fill=S1BgCol, minimum width=27mm},
    ablk/.style={blk, draw=BranchABdr, fill=BranchABg, minimum width=24mm},
    bblk/.style={blk, draw=BranchBBdr, fill=BranchBBg, minimum width=24mm},
    fblk/.style={blk, draw=FuseBdr, fill=FuseBg, minimum width=30mm},
    embnode/.style={
        draw, rounded corners=3pt,
        inner sep=3pt, minimum height=9mm, minimum width=18mm,
        line width=0.55pt
    },
    auxblk/.style={
        draw=AuxCol, fill=AuxBg, dashed,
        rounded corners=2pt,
        inner sep=3pt, minimum height=7mm, minimum width=22mm,
        line width=0.45pt, font=\footnotesize\sffamily
    },
    splitnode/.style={
        circle, draw=IOCol, fill=white,
        inner sep=0pt, minimum size=4pt,
        line width=0.5pt
    },
    fl/.style={->, line width=0.5pt, draw=black!75},
    flA/.style={->, line width=0.5pt, draw=BranchACol},
    flB/.style={->, line width=0.5pt, draw=BranchBCol},
    flF/.style={->, line width=0.6pt, draw=FuseCol},
    auxfl/.style={->, dashed, line width=0.4pt, draw=AuxCol},
    tlbl/.style={font=\scriptsize\itshape\sffamily, text=ShapeCol, inner sep=1pt},
    losslbl/.style={font=\scriptsize\sffamily, text=S1Blue, inner sep=1pt},
    auxlbl/.style={font=\scriptsize\sffamily, text=AuxCol, inner sep=1pt},
]

%% ═══════════════════════════════════════════════════════════════════════════
%%   [Input]──[Denoiser]──[Denoised]──●──┬──[Reshape]──[ResNet]──[ProjH]──[e_eeg]──┐
%%                                       │                                          ├──[Fusion]──[f]──[Loss]
%%                                       └──[STFT]──[ConvStem]──[SpecMamba]──[e_spec]┘
%% ═══════════════════════════════════════════════════════════════════════════

%% ── Stage I: input → denoiser → denoised ──
\node[ionode] (input) {%
    \textbf{Raw EEG}\\
    $\mathbb{R}^{4\times 800}$%
};

\node[s1blk, right=10mm of input] (denoiser) {%
    \textcolor{S1Blue}{\textbf{WaveNet + Mamba}}\\
    \footnotesize 3 blocks $\times$ 4 dil.\ layers\\
    \footnotesize + Mamba SSM (mid)%
};
\node[losslbl, below=1.2mm of denoiser] (sisnrlbl) {\textit{SI-SNR loss}};

\node[ionode, right=10mm of denoiser] (denoised) {%
    \textbf{Denoised} $\hat{x}$\\
    $\mathbb{R}^{4\times 800}$%
};

\node[splitnode, right=11mm of denoised] (split) {};

%% ── Branch A (upper): ResNet EEG embedder ──
\node[ablk, above right=7mm and 7mm of split] (reshape) {%
    \textcolor{BranchACol}{\textbf{Reshape}}\\
    \footnotesize $(4,800)\to(4,25,32)$%
};

\node[ablk, right=11mm of reshape] (resnet) {%
    \textcolor{BranchACol}{\textbf{ResNet-18/34}}\\
    \footnotesize $3{\times}3$ stride-1 stem\\
    \footnotesize (no maxpool)%
};

\node[ablk, right=11mm of resnet] (projhead) {%
    \textcolor{BranchACol}{\textbf{Proj.\ Head}}\\
    \footnotesize Lin$\to$ReLU$\to$Drop\\
    \footnotesize $\to$Lin$\to$BN%
};

\node[embnode, draw=BranchABdr, fill=BranchABg, right=11mm of projhead] (eegemb) {%
    $e_{\mathrm{eeg}}\!\in\!\mathbb{R}^{128}$\\
    \footnotesize L2-norm%
};

%% ── Branch B (lower): Spectrogram Mamba ──
\node[bblk, below right=7mm and 7mm of split] (stft) {%
    \textcolor{BranchBCol}{\textbf{STFT}}\\
    \footnotesize $n_{\mathrm{fft}}{=}128$\\
    \footnotesize hop$=64$%
};

\node[bblk, right=11mm of stft] (convstem) {%
    \textcolor{BranchBCol}{\textbf{Conv Stem}}\\
    \footnotesize $4{\to}32{\to}64$ ch.\\
    \footnotesize $3{\times}3$, BN-ReLU%
};

\node[bblk, right=11mm of convstem] (specmamba) {%
    \textcolor{BranchBCol}{\textbf{Spec.\ Mamba}}\\
    \footnotesize freq sweep\\
    \footnotesize $+$ time sweep%
};

\node[embnode, draw=BranchBBdr, fill=BranchBBg, right=11mm of specmamba] (specemb) {%
    $e_{\mathrm{spec}}\!\in\!\mathbb{R}^{128}$\\
    \footnotesize L2-norm%
};

%% ── Fusion (vertically centred between embeddings) ──
\node[fblk] at ($(eegemb.east)!0.5!(specemb.east) + (14mm,0)$) (fusion) {%
    \textcolor{FuseCol}{\textbf{Gate Fusion}}\\
    \footnotesize learned 128-d gate\\
    \footnotesize $+$ residual MLP%
};

\node[embnode, draw=FuseBdr, fill=FuseBg, right=8mm of fusion] (fout) {%
    $f\!\in\!\mathbb{R}^{128}$\\
    \footnotesize L2-norm%
};

\node[fblk, right=11mm of fout, minimum width=24mm] (mainloss) {%
    \textcolor{FuseCol}{\textbf{Metric Loss}}\\
    \footnotesize ArcFace / Multi-Sim%
};

%% ── Auxiliary loss boxes (sit just outside the branches, kept compact) ──
\node[auxblk, above=4mm of eegemb] (auxeeg) {%
    \textcolor{AuxCol}{$\mathcal{L}_{\mathrm{metric}}(e_{\mathrm{eeg}})$, $\lambda_e{=}0.3$}%
};

\node[auxblk, below=4mm of specemb] (auxspec) {%
    \textcolor{AuxCol}{$\mathcal{L}_{\mathrm{metric}}(e_{\mathrm{spec}})$, $\lambda_s{=}0.2$}%
};

%% ═══════════════════════════════════════════════════════════════════════════
%% ARROWS
%% ═══════════════════════════════════════════════════════════════════════════

\draw[fl] (input) -- node[tlbl, above] {$\mathbb{R}^{4\times800}$} (denoiser);
\draw[fl] (denoiser) -- node[tlbl, above] {$\mathbb{R}^{4\times800}$} (denoised);
\draw[fl] (denoised) -- (split);

%% Split → Branch A
\draw[flA] (split) |- (reshape.west);
\draw[flA] (reshape) -- (resnet);
\draw[flA] (resnet)  -- (projhead);
\draw[flA] (projhead) -- node[tlbl, above] {$\mathbb{R}^{256}$} (eegemb);

%% Split → Branch B
\draw[flB] (split) |- (stft.west);
\draw[flB] (stft)     -- node[tlbl, above] {$4{\times}65{\times}13$} (convstem);
\draw[flB] (convstem) -- (specmamba);
\draw[flB] (specmamba) -- node[tlbl, above] {$\mathbb{R}^{256}$} (specemb);

%% Both embeddings → fusion (enter from top/bottom to avoid crossing through box)
\draw[flF] (eegemb.east)  -| (fusion.north);
\draw[flF] (specemb.east) -| (fusion.south);

%% Fusion → output → loss
\draw[flF] (fusion) -- node[tlbl, above] {$\mathbb{R}^{128}$} (fout);
\draw[flF] (fout)   -- (mainloss);

%% Auxiliary loss arrows
\draw[auxfl] (eegemb.north)  -- (auxeeg.south);
\draw[auxfl] (specemb.south) -- (auxspec.north);

%% ═══════════════════════════════════════════════════════════════════════════
%% BACKGROUND STAGE BOXES (aux losses excluded — they sit visually outside)
%% ═══════════════════════════════════════════════════════════════════════════
\begin{pgfonlayer}{background}

    \node[
        draw=S1Bdr, dashed, rounded corners=4pt,
        fill=S1BgCol, fill opacity=0.30,
        line width=0.5pt,
        fit=(input)(denoiser)(denoised)(sisnrlbl),
        inner sep=4pt
    ] (stage1box) {};

    \node[
        draw=FuseBdr, dashed, rounded corners=4pt,
        fill=FuseBg, fill opacity=0.20,
        line width=0.5pt,
        fit=(reshape)(resnet)(projhead)(eegemb)
            (stft)(convstem)(specmamba)(specemb)
            (fusion)(fout)(mainloss),
        inner sep=4pt
    ] (stage2box) {};

\end{pgfonlayer}

%% Stage labels — placed inside top-left of each box, small caps
\node[
    font=\footnotesize\bfseries\sffamily,
    text=S1Blue,
    above=1pt of stage1box.north west,
    anchor=south west
] {\textsc{Stage I}: Denoising};

\node[
    font=\footnotesize\bfseries\sffamily,
    text=FuseCol,
    above=1pt of stage2box.north west,
    anchor=south west
] {\textsc{Stage II}: Multimodal Embedding $+$ Fusion};

\end{tikzpicture}
}
\caption{V4 multimodal architecture. The WaveNet+Mamba denoiser feeds two
parallel branches: a ResNet metric embedder on the denoised waveform and
a spectrogram Mamba branch with frequency- and time-axis sweeps. A
gated self-attention fusion module fuses the two L2-normalized
embeddings into a single 128-d biometric embedding. Auxiliary metric losses
on each branch (dashed) prevent fusion collapse during joint training.}
\label{fig:v4-arch}
\end{figure*}

\section{Experimental Setup}
\label{sec:setup}
\textbf{Dataset.} We use the publicly available PhysioNet
\emph{Auditory-Evoked-Potential EEG-Biometric} corpus~\cite{alzahab2024aep},
recording 20 subjects across 12 sessions of approximately two minutes
each under both resting-state and auditory-stimulation conditions, with
in-ear and bone-conduction headphones. Auditory evoked potentials are
the cortical response to a controlled acoustic stimulus; their
latency- and amplitude-domain morphology can carry subject-specific
structure that originates in individual differences
of cortical anatomy and auditory processing pathways, which makes them
attractive as an internal, replay-resistant biometric cue.
Recordings are segmented into fixed-length 4-channel epochs of
$T=800$ samples.

\textbf{Evaluation protocol.} We inherit the subject partitioning and
open-set scoring rule of~\cite{vu2026aisei} and refer to it as the
\emph{AEP-Hybrid protocol} for brevity below. Concretely, 4 of the
20 subjects ($\mathcal{S}_{\mathrm{holdout}}=\{2,5,7,12\}$) are
strictly held out as an unknown-impostor pool, and the remaining 16
known subjects are split per-subject into train/validation/test along
sessions and epochs. Two evaluation tasks share this split.
\emph{(i) Closed-set retrieval} reports P@1 and P@5 on the test split
of the 16 known subjects (session- and epoch-disjoint, identity space
seen at training).
\emph{(ii) Held-out-subject open-set scoring} forms a per-subject
prototype as the mean of training-set fused embeddings, scores each
held-out epoch by its minimum Euclidean distance to any prototype, and
fixes the operating threshold at the 95th percentile of the same
min-distance on validation. Held-out-subject AUROC under this scoring
stays near chance ($\approx 0.50$) and is reported in the discussion;
the AUROC and EER columns in Table~\ref{tab:v4} are closed-set
pairwise verification scores on the same test pool, characterizing
within-subject session variability rather than held-out-subject
novelty. We repeat each configuration with three random seeds
(PyTorch, NumPy, and Python \texttt{random} seeds $\{0,1,2\}$) and
report mean$\pm$standard deviation.

\textbf{Noise families.} We synthesize three noise types at deployment:
white Gaussian, $50$\,Hz sinusoidal power-line, and broadband
$20$--$80$\,Hz EMG-like noise.

\textbf{Hyperparameters.} The auxiliary-loss weights are fixed at
$\lambda_{e}=0.3$ and $\lambda_{s}=0.2$ and the class-balanced sampler
parameter at $m=4$. These operating points lie on the plateau of
single-seed sensitivity sweeps over $\lambda_{e}\!\times\!\lambda_{s}$
and $m\in\{2,4,8,16\}$ on the R34+ArcFace + power-line configuration;
the full grids are deferred to the journal version, while the headline
three-seed runs in Tables~\ref{tab:cross} and~\ref{tab:v4} use the
fixed operating point above.

\textbf{Architectures.} We compare four variants under identical
data, splits, seeds, and training schedules:
\begin{itemize}
\item \textbf{V1}: WaveNet denoiser only (no Mamba) + EEG-only metric
embedder. This is the prior pipeline of~\cite{vu2026aisei}
re-evaluated under the AEP-Hybrid protocol.
\item \textbf{V2}: V1 with the midpoint Mamba block in the
denoiser.
\item \textbf{V3}: V2 with a tuned Mamba preset, run as a directional,
single-seed, different-hardware quick run (NVIDIA H100 with
\texttt{torch.compile}); not used for the controlled three-seed
claims.
\item \textbf{V4}: full multimodal model proposed in this paper:
WaveNet+Mamba denoiser, ResNet EEG branch, Mamba spectrogram branch, and
gated self-attention fusion with auxiliary branch losses.
\end{itemize}

\textbf{Implementation.} All three-seed V1, V2, and V4 experiments run
on a single NVIDIA RTX~5090 with batch size~$256$, $8$ data-loader
workers, mixed precision, and the \texttt{denoised} spectrogram source
for V4. The V4 model totals $22.8$\,M parameters with end-to-end
inference latency of $186$--$197$\,ms per $4$\,s epoch on an RTX~5090
(ResNet-34+ArcFace head), well below the $4$\,s epoch duration and
therefore real-time-compatible; the V1 baseline reaches $\sim\!20$\,ms
on an RTX~3060 with comparable parameter count
($\sim\!23.4$\,M)~\cite{vu2026aisei}, so the added Mamba blocks,
spectrogram branch, and gated fusion module account for the latency
delta on top of the V1 backbone. V3 was a single-seed quick run on an
NVIDIA H100 with \texttt{torch.compile}, reported as directional only.

\textbf{Metrics.} We distinguish a \emph{score} (a cosine similarity
over a single pair of embeddings) from a \emph{metric} (an
aggregation over the score population). For closed-set retrieval we
report Precision@1 (P@1) and Precision@5 (P@5): P@5 is the fraction
of the top-5 retrieved gallery epochs (excluding the query itself)
that share the query identity, averaged across queries. For signal
quality we report SI-SNR in decibels. For closed-set pairwise
verification we report AUROC and equal-error rate (EER):
AUROC is the area under the receiver-operating-characteristic curve
constructed by sweeping a threshold over the cosine-similarity score,
and EER is the threshold at which false-accept and false-reject
rates coincide. Both are computed over all test-pool embedding pairs
(Table~\ref{tab:v4}); held-out-subject open-set scoring is described
separately in the protocol paragraph.

\section{Results and Discussion}
\subsection{Cross-Version Best P@1}
Table~\ref{tab:cross} summarizes the best Precision@1 of each version,
selected per noise family across the three metric heads. The proposed
V4 model is the strongest configuration on every noise type, with the
clearest margins on power-line and EMG. V3 reaches a competitive
power-line score (its tuned Mamba preset narrows the spectral residual
the $50$\,Hz line leaves behind), but its Gaussian and EMG numbers
trail V1, V2, and V4: with a single seed and a different hardware
target (H100 + \texttt{torch.compile}) the run is not directly
comparable to the three-seed RTX~5090 numbers and is reported only as
a directional reference.

\begin{table}[t]
\centering
\caption{Best Precision@1 (P@1) per version and noise type. Higher is
better. V3 is a single-seed quick run.}
\label{tab:cross}
\begin{threeparttable}
\begin{tabular}{l c c c c}
\toprule
\textbf{Noise} & \textbf{V1} & \textbf{V2} & \textbf{V3}\tnote{$\dagger$}
& \textbf{V4} \\
\midrule
Gaussian   & 0.822 & 0.798 & 0.749 & \textbf{0.824} \\
Power-line & 0.860 & 0.858 & 0.869 & \textbf{0.877} \\
EMG        & 0.824 & 0.811 & 0.758 & \textbf{0.838} \\
\bottomrule
\end{tabular}
\begin{tablenotes}\footnotesize
\item[$\dagger$] Single-seed H100 quick run; reported for reference only.
\end{tablenotes}
\end{threeparttable}
\end{table}

% Auto-generated by experiments/analysis/build_priorwork_table.py.
% Re-run that script to regenerate; do not edit by hand.
\begin{table}[t]
\centering
\caption{Cited prior-work numbers on EEG-based identification. Each row uses its source paper's protocol; no baseline is re-implemented under the AEP-Hybrid protocol of this paper. V4 results are in Table~\ref{tab:cross}.}
\label{tab:priorwork}
\begin{threeparttable}
\scriptsize
\setlength{\tabcolsep}{3pt}
\begin{tabular}{l l c c c}
\toprule
\textbf{Method} & \textbf{Dataset} & \textbf{Subj.} & \textbf{Metric} & \textbf{Score} \\
\midrule
MindID~\cite{zhang2017mindid}\tnote{1} & EEG-ID & 8 & Acc. & 0.984 \\
EEG-Triplet~\cite{yang2022eeg}\tnote{2} & PhysioNet MI & 109 & Acc. & 0.961 \\
BrainNet~\cite{fallahi2023brainnet}\tnote{3} & ERP-CORE & 40 & EER & 0.027 \\
AISEI v0~\cite{vu2026aisei}\tnote{4} & PhysioNet AEP & 20 & P@1 & 0.860 \\
\bottomrule
\end{tabular}
\begin{tablenotes}\scriptsize
\item[1] Within-session resting-state EEG; attention-RNN; no synthetic noise injection.
\item[2] Motor-imagery task; triplet loss on raw signals; session-mixed split.
\item[3] Siamese network on ERP epochs; cross-session evaluation; open-set verification on enrolled subjects only.
\item[4] Same dataset and same AEP-Hybrid protocol as this work; WaveNet-only denoiser; power-line noise.
\end{tablenotes}
\end{threeparttable}
\end{table}


\subsection{Detailed V4 Results}
Table~\ref{tab:v4} reports the full V4 metric grid for the three
metric heads. ResNet-34 with ArcFace is consistently the strongest head
in P@1, while ResNet-18 with multi-similarity is competitive on
power-line. The denoiser reaches its highest SI-SNR in the
power-line condition (the $50$\,Hz sinusoid is structurally easy to
remove); Gaussian and EMG SI-SNR are lower but still in the regime
expected for non-stationary residuals.

\begin{table}[t]
\centering
\caption{V4 multimodal results, mean$\pm$std over three seeds. The
ResNet-34+ArcFace head is the strongest by P@1 on all noise families.}
\label{tab:v4}
\begin{threeparttable}
\scriptsize
\setlength{\tabcolsep}{2.2pt}
\begin{tabular}{l l c c c c c}
\toprule
\textbf{Noise} & \textbf{Head} & \textbf{P@1} & \textbf{P@5}
& \textbf{SI-SNR} & \textbf{AUROC} & \textbf{EER} \\
\midrule
\multirow{3}{*}{Gauss.}
 & R34+MS  & .764$\pm$.050 & .743$\pm$.053 & 12.15$\pm$.26 & .443$\pm$.021 & .376$\pm$.072 \\
 & R18+MS  & .782$\pm$.028 & .755$\pm$.035 & 12.15$\pm$.26 & .441$\pm$.036 & .369$\pm$.046 \\
 & R34+Arc & \textbf{.824$\pm$.036} & \textbf{.803$\pm$.040} & 12.15$\pm$.26 & .479$\pm$.095 & .356$\pm$.059 \\
\midrule
\multirow{3}{*}{P-line}
 & R34+MS  & .819$\pm$.040 & .794$\pm$.049 & 32.38$\pm$1.06 & .533$\pm$.051 & .408$\pm$.024 \\
 & R18+MS  & .831$\pm$.047 & .806$\pm$.052 & 32.38$\pm$1.06 & .450$\pm$.015 & .371$\pm$.021 \\
 & R34+Arc & \textbf{.877$\pm$.006} & \textbf{.855$\pm$.005} & 32.38$\pm$1.06 & .501$\pm$.044 & .360$\pm$.039 \\
\midrule
\multirow{3}{*}{EMG}
 & R34+MS  & .764$\pm$.080 & .747$\pm$.075 & 13.83$\pm$.35 & .439$\pm$.065 & .389$\pm$.039 \\
 & R18+MS  & .793$\pm$.067 & .775$\pm$.069 & 13.83$\pm$.35 & .422$\pm$.038 & .356$\pm$.066 \\
 & R34+Arc & \textbf{.838$\pm$.057} & \textbf{.827$\pm$.065} & 13.83$\pm$.35 & .514$\pm$.082 & .350$\pm$.058 \\
\bottomrule
\end{tabular}
\begin{tablenotes}\footnotesize
\item MS: multi-similarity loss; Arc: ArcFace; R18/R34: ResNet-18/34.
\end{tablenotes}
\end{threeparttable}
\end{table}

\subsection{Ablation: Mamba and Multimodal Fusion}
Table~\ref{tab:ablation} isolates the contribution of each architectural
ingredient on the strongest head (ResNet-34 + ArcFace) using P@1 as the
discriminating metric. Inserting a single Mamba block into the WaveNet
denoiser (V1$\to$V2) leaves P@1 essentially unchanged and even slightly
hurts it, suggesting that swapping a single fixed-receptive-field
convolutional layer for a state-space block does not, on its own,
re-shape the embedding geometry. Adding the spectrogram branch and
gated self-attention fusion (V2$\to$V4) is what unlocks retrieval gains of
$+2.6$, $+1.9$, and $+2.7$ percentage points respectively over V2, and
$+0.2$, $+1.7$, and $+1.4$ over V1. We do not attempt to attribute
SI-SNR deltas to the Stage~II additions: each version retrains its own
Stage~I denoiser independently, so cross-version SI-SNR differences
reflect independent training runs of nominally the same architecture
rather than any contribution from the spectrogram branch or fusion
module. Per-version Stage~I SI-SNR (mean across seeds, ResNet-34+ArcFace
configuration) is V1: $10.64/19.98/11.65$\,dB; V2:
$10.65/19.61/11.67$\,dB; V4: $12.15/32.38/13.83$\,dB on Gaussian /
power-line / EMG, and is intended only as descriptive Stage~I quality.

\begin{table}[t]
\centering
\caption{Architectural ablation on ResNet-34+ArcFace by P@1.
$\Delta$ rows are V4 minus V1. Stage~I is retrained independently per
version, so SI-SNR is not used here as an ablation metric.}
\label{tab:ablation}
\begin{tabular}{l l c}
\toprule
\textbf{Noise} & \textbf{Variant} & \textbf{P@1} \\
\midrule
\multirow{4}{*}{Gaussian}
 & V1: WaveNet  & 0.822 \\
 & V2: + Mamba  & 0.798 \\
 & V4: + Spec + Fusion & \textbf{0.824} \\
 & $\Delta$ V4$-$V1 & $+0.002$ \\
\midrule
\multirow{4}{*}{Power-line}
 & V1: WaveNet  & 0.860 \\
 & V2: + Mamba  & 0.858 \\
 & V4: + Spec + Fusion & \textbf{0.877} \\
 & $\Delta$ V4$-$V1 & $+0.017$ \\
\midrule
\multirow{4}{*}{EMG}
 & V1: WaveNet  & 0.824 \\
 & V2: + Mamba  & 0.811 \\
 & V4: + Spec + Fusion & \textbf{0.838} \\
 & $\Delta$ V4$-$V1 & $+0.014$ \\
\bottomrule
\end{tabular}
\end{table}

\subsection{Discussion}
\textbf{Where multimodal helps.} The largest mean-P@1 jumps occur on
power-line noise, consistent with the spectrogram localizing the
$50$\,Hz line and the gate suppressing it before fusion. EMG is
intermediate (broadband but high-frequency-tilted); Gaussian gains are
smallest, consistent with white noise being structureless in both
views.

\textbf{Effect sizes and the Mamba-alone variant.} V4-over-V1 and
V2-vs-V1 mean-P@1 deltas both fall within across-seed standard
deviation, so we report them as magnitude observations rather than
tested claims (see Limitations). The implication is that the gains
should be read as a property of the full multimodal stack — denoiser,
spectrogram branch, and gated fusion together — rather than an
isolated Stage-I Mamba effect.

\textbf{Verification gap.} Held-out-subject AUROC of the centroid
min-distance score stays close to chance ($\approx 0.50$) across all
V4 configurations, confirming that threshold-free verification of
unseen identities remains the open problem. The AUROC ($0.42$--$0.53$)
and EER ($0.35$--$0.41$) columns in Table~\ref{tab:v4} are computed on
the closed-set test pool (session- and epoch-disjoint, all subjects
seen during training) and indicate within-subject session variability,
not novel-identity verification; calibration is left to future work.

\textbf{Limitations and future work.} Each limitation below is paired
with the planned remedy. (i) The dataset is small (20 subjects, 4
channels), so absolute deltas are within across-seed noise; we plan a
multi-cohort recording with $\geq 100$ subjects to support paired
significance at the per-noise level. (ii) Noise is synthetic by
controlled injection; we plan a live-environment recording session
with hardware artefacts (movement, blink, electrode drift) to retire
the synthetic-only caveat. (iii) Held-out-subject AUROC is near
chance; we plan verification calibration via Platt scaling on a
held-out validation cohort and a per-subject score-normalisation
study. (iv) Same-protocol re-implementation of external EEG-biometric
baselines (EEGNet, ChronoNet, ICA+SVM) and a cross-dataset transfer
evaluation are out of scope here and are scheduled for the journal
version of this work; Table~\ref{tab:priorwork} collects cited
numbers under their own protocols as an informational, not controlled,
positioning of the V4 result.

\section{Conclusion}
We presented a multimodal Mamba-augmented framework for noise-robust
EEG biometric identification extending~\cite{vu2026aisei}: V4 combines
a WaveNet+Mamba denoiser, a spectrogram Mamba branch, and a gated
self-attention fusion module trained with auxiliary branch losses.
Under the AEP-Hybrid protocol (Sec.~\ref{sec:setup}) on the PhysioNet
AEP dataset, V4 attains the highest mean Precision@1 on all three
noise families ($82.4\%$/$87.7\%$/$83.8\%$); held-out-subject AUROC
$\approx 0.50$ motivates calibration as future work.

\AtNextBibliography{\footnotesize}
\setlength{\bibitemsep}{0pt plus 0.3ex}
\printbibliography

\end{document}
